% Bagian: computability
% Bab: computability-theory
% Subbagian: applications-fixed-point

\documentclass[../../../include/open-logic-section]{subfiles}

\begin{document}

\olfileid[id]{cmp}{thy}{apf}
\olsection{Penerapan Teorema Titik Tetap}

Teorema titik tetap pada dasarnya memungkinkan kita mendefinisikan fungsi
dapat dihitung parsial berdasarkan indeksnya. Sebagai contoh, kita dapat
menemukan indeks~$e$ sedemikian sehingga bagi setiap $y$,
\[
\cfind{e}(y) = e+y.
\]
Sebagai contoh lain, bukti teorema titik tetap dapat digunakan untuk merancang
program dalam Java atau C++ yang mencetak dirinya sendiri.

Ingat bahwa jika bagi setiap $e$ kita menetapkan $W_e$ sebagai domain
$\cfind{e}$, maka barisan $W_0$, $W_1$, $W_2$,~\dots mengenumerasi
himpunan-himpunan yang terenumerasi secara komputabel. Beberapa himpunan ini
dapat dihitung. Kita dapat bertanya apakah terdapat algoritma yang menerima
nilai~$x$ sebagai masukan dan, jika $W_x$ kebetulan dapat dihitung,
mengembalikan indeks bagi fungsi karakteristiknya. Jawabannya ``tidak'';
algoritma semacam itu tidak ada:

\begin{thm}
Tidak ada fungsi dapat dihitung parsial $f$ dengan sifat berikut: setiap kali
$W_e$ dapat dihitung, $f(e)$ terdefinisi dan $\cfind{f(e)}$ merupakan fungsi
karakteristiknya.
\end{thm}

\begin{proof}
Andaikan, demi memperoleh kontradiksi, bahwa $f$ merupakan fungsi dapat
dihitung parsial dengan sifat dalam pernyataan teorema. Dengan menggunakan
pernyataan pertama teorema titik tetap pada fungsi~$g$ yang didefinisikan di
bawah, kita dapat menemukan indeks~$e$ sedemikian sehingga
\[
\cfind{e}(y) \simeq
\begin{cases}
0 & \text{jika $y=0$ dan $\cfind{f(e)}(0) \fdefined = 0$} \\
\text{tidak terdefinisi} & \text{selain itu.}
\end{cases}
\]
Fungsi yang digunakan dalam penerapan teorema titik tetap tersebut ialah
\[
g(x,y) \simeq
\begin{cases}
0 & \text{jika $y=0$ dan $\cfind{f(x)}(0) \fdefined = 0$} \\
\text{tidak terdefinisi} & \text{selain itu.}
\end{cases}
\]
Secara informal, kita dapat melihat bahwa $g$ dapat dihitung parsial sebagai
berikut: pada masukan $x$ dan $y$, algoritma mula-mula memeriksa apakah $y$
sama dengan~$0$. Jika ya, algoritma menghitung $f(x)$, lalu menggunakan mesin
universal untuk menghitung $\cfind{f(x)}(0)$. Jika komputasi terakhir ini
berhenti dan mengembalikan~$0$, algoritma mengembalikan~$0$; jika tidak,
algoritma tidak berhenti.

Sekarang, perhatikan bahwa $W_e$ selalu sama dengan $\{0\}$ atau
$\emptyset$, sehingga $W_e$ dapat dihitung. Sifat yang diandaikan bagi~$f$
karena itu mengharuskan $f(e)$ terdefinisi dan $\cfind{f(e)}$ menjadi fungsi
karakteristik~$W_e$. Namun, jika $\cfind{f(e)}(0)$ terdefinisi dan sama
dengan~$0$, maka $\cfind{e}(y)$ terdefinisi tepat ketika $y=0$, sehingga
$W_e=\{0\}$ dan $\cfind{f(e)}$ memberikan nilai yang salah pada~$0$. Jika
$\cfind{f(e)}(0)$ tidak terdefinisi, atau terdefinisi tetapi tidak sama
dengan~$0$, maka $W_e=\emptyset$ dan sekali lagi $\cfind{f(e)}$ bukan fungsi
karakteristik~$W_e$. Kedua kemungkinan menghasilkan kontradiksi.
\end{proof}

\end{document}
